% !PRINTED 321
\weberpage{321}
% [page alignment mirrored 2026-08-07] The three blocks that opened this fragment (tail of
% XIII, connecting sentence, theorem XIV) are printed on page 320 and were moved to i0345,
% following the corrected German. Page 321 begins directly with "Diese Sätze werden…".
These theorems are substantially generalised if one applies a linear
transformation to them\footnote{\gesperrt{Laguerre}, Nouvelles
Annales de Math\'ematiques, 2.\,S\'er., Vol.\,19, 1880.}.

This is more easily carried out if one uses the homogeneous notation,
that is, understands by $f(x,y)$ an integral homogeneous function
of degree $n$. If, for an arbitrary value $\xi:\eta$, one forms the polars
of this function (\S.\,60):
\[
P_1(x,\xi) = \frac{1}{n}\{\xi f'(x) + \eta f'(y)\}
\]
\weq{3}{P_2(x,\xi) = \frac{1}{n(n-1)}\{\xi^2 f''(x,x) + 2\xi\eta f''(x,y) + \eta^2 f''(y,y)\}}
\[
P_3(x,\xi) = \frac{1}{n(n-1)(n-2)}\{\xi^3 f'''(x,x,x) + 3\xi^2\eta f'''(x,x,y) + 3\xi\eta^2 f'''(x,y,y) + \eta^3 f'''(y,y,y)\}
\]
\[
\cdots\cdots\cdots\cdots\cdots\cdots\cdots\cdots\cdots\cdots\cdots\cdots,
\]
then these functions remain unchanged if one transforms $\xi, \eta$ and $x, y$
simultaneously by a linear transformation and puts for $f$ and
its derivatives the corresponding transformed functions. With
respect to the reality of the roots, nothing is changed by a real linear
transformation. But now, if $\xi, \eta$ are real,
the real linear substitution can be so determined that the transformed
values become $\xi' = 1$, $\eta' = 0$, and thereby, if one also sets
$y = 1$, the polars pass into the derivatives of the function $f(x)$ with respect to $x$
(apart from numerical factors).

\gesperrt{From this it follows that in the theorems XIII, XIV the derivatives $f'(x), f''(x)
\dots$ can be replaced by the polars $P_1(x,\xi), P_2(x,\xi) \dots$ for arbitrary
real $\xi, \eta$.}

We make an application of this to the equation which one obtains by
setting the $(n-2)$th polar equal to zero:
\weq{4}{x^2 f''(\xi,\xi) + 2xy f''(\xi,\eta) + y^2 f''(\eta,\eta) = 0.}
If $f(x,y) = 0$ has, as we assume, only real roots, then the
quadratic equation (4) must also have real roots, i.\,e. the
Hesse determinant
\weq{5}{f''(\xi,\xi) f''(\eta,\eta) - f''(\xi,\eta)^2}
must be negative for all real $\xi, \eta$. It can only van-
%CONT>>
